%% ACM SIGCHI two-column format.
%% Compile locally: pdflatex main && bibtex main && pdflatex main && pdflatex main
%% Or paste section content into the Overleaf acmart template:
%% https://www.overleaf.com/latex/templates/association-for-computing-machinery-acm-sig-proceedings-template/bmvfhcdnxfty

\documentclass[sigconf,screen]{acmart}

\acmConference[CHI '26]{ACM CHI Conference on Human Factors in Computing Systems}{April 26--May 1, 2026}{Yokohama, Japan}
\acmDOI{}
\acmISBN{}
\setcopyright{rightsretained}
\settopmatter{printacmref=false}

\usepackage{booktabs}
\usepackage{microtype}
\usepackage{graphicx}
\usepackage{xcolor}

\begin{document}

\title{Easy\_SLM: A Scaffolded Interface for Structured\\
       SLM--LLM Comparison on Everyday Tasks}

\author{Jun Hou}
\affiliation{%
  \institution{[Institution]}
  \country{[Country]}}
\email{kaitlin.jun@gmail.com}

\begin{abstract}
Everyday users increasingly face a choice between private, offline-capable
small language models (SLMs) and more powerful remote large language models
(LLMs)---yet fair comparisons are confounded by prompting-skill variance and
interface differences.
We first conducted a formative Reddit study collecting 3,292 posts and
comments from SLM and LLM communities (March--May 2025).
GPT-4o annotation followed by canonical clustering identified 29 task groups,
19 of which (65\%) were shared across both streams, with creative writing,
summarization, scripting, and planning most frequent.
Grounded in these findings and the Gulf of Envisioning framework, we
designed \textit{Easy\_SLM}: an interface with four design goals---paired
SLM--LLM execution with objective scoring, structured goal specification,
prompt scaffolding, and guided model selection.
We describe a within-subjects user study ($N$~=~[\textit{N}]) in which
participants completed writing, summarizing, and planning tasks in both
conditions.
Our results quantify where SLM outputs are ``good enough,'' how scaffolding
reduces prompting-skill variance, and how interface affordances trade off
workload and output quality between the two model classes.
\end{abstract}

\keywords{small language models; Gulf of Envisioning; prompt scaffolding;
SLM--LLM comparison; sustainable AI; everyday tasks}

\maketitle

% =============================================================================
\section{Introduction}
% =============================================================================

Everyday users increasingly interact with AI assistants for writing,
summarization, planning, and information search.
Two broad classes of model are now practically available: \textit{large language
models} (LLMs), accessed via cloud APIs and characterized by high accuracy and
robustness, and \textit{small language models} (SLMs), which run locally on
commodity hardware, offer stronger privacy guarantees, and operate without
internet connectivity.
The pragmatic question---\textit{for a given daily task, is an SLM good enough?}
---has real stakes for sustainable computing and equitable access, yet it has no
clean empirical answer.

The difficulty is not purely technical.
Observed differences in output quality between an SLM and an LLM are
\textit{entangled} with at least two confounds: (i)~\textit{prompting-skill
variance}, because users write different prompts for the two systems and SLMs are
more brittle to underspecified input~\cite{zamfirescu2023johnny,
subramonyam2024gulf}; and (ii)~\textit{interface differences}, because most
SLM deployments offer a bare chat box while commercial LLM products embed
suggestion, feedback, and correction affordances.
Letting users ``freestyle'' prompts in different tools and comparing outputs
conflates model capability with user skill and UI support---making it impossible
to attribute gaps to the model.

This paper addresses both confounds through a single shared interface,
\textit{Easy\_SLM}, that routes the same user specification to either an SLM
(via effGen~\cite{effgen}) or a cloud LLM (OpenAI / Anthropic), and that
provides \textit{light scaffolding} to control Gulf-of-Envisioning
variance~\cite{subramonyam2024gulf} in both conditions.
Our work is motivated by three concerns.
\textbf{Sustainability:} ``Use a smaller model'' is not actionable without
knowing when an SLM is genuinely sufficient~\cite{crawford2024rematerializing}.
\textbf{Resilient access:} Cloud dependency creates brittleness; the 2024
CrowdStrike outage disrupted AI-assisted care in hundreds of hospitals, and
offline-capable SLMs offer essential fallback~\cite{tully2025crowdstrike}.
\textbf{Democratization:} Novice users face a Gulf of Envisioning with SLMs
because these models are less forgiving of vague intent; scaffolding shifts
``prompting as skill'' to ``prompting as structure''~\cite{zamfirescu2023johnny}.

We address four research questions:
\textbf{RQ1} (Capability gap): How does SLM performance compare to LLM
performance on daily tasks under objective metrics (ROUGE, F1, rule-based
pass rates)?
\textbf{RQ2} (``Good enough'' boundary): For which tasks can SLM outputs be
considered ``good enough'' based on combined objective scores and user
preference?
\textbf{RQ3} (Prompting confound): To what extent does a standardized
prompting interface reduce prompting-skill variance in SLM--LLM comparisons?
\textbf{RQ4} (Affordances): How do instruction decomposition and
constraint-checking affordances affect output quality, format violations, and
perceived workload across conditions?

Our contributions are: (1)~a formative Reddit study identifying 19 everyday
tasks shared across SLM and LLM user communities; (2)~the Easy\_SLM system
with four design goals operationalizing fair paired comparison; and (3)~user
study evidence on capability gaps, ``good enough'' task boundaries, and
scaffold effectiveness.

% =============================================================================
\section{Related Work}
% =============================================================================

\subsection{Gulf of Envisioning in LLM Interaction}

Subramonyam et al.\ identify a \textit{Gulf of Envisioning} in LLM-based
interaction: users struggle to form mental models of what a model can or will
produce, making it hard to specify intent precisely enough for reliable
output~\cite{subramonyam2024gulf}.
This gap is especially severe for SLMs, which lack the instruction-following
robustness of frontier models~\cite{zamfirescu2023johnny}.
Zamfirescu-Pereira et al.\ show that non-expert users systematically fail to
craft effective prompts even for LLMs, underspecifying constraints, omitting
output-format expectations, and iterating without a principled
strategy~\cite{zamfirescu2023johnny}.
Easy\_SLM addresses the Gulf of Envisioning specifically in the SLM context by
externalizing goal structure, decomposing prompts into steps, and capturing
audience and format constraints as first-class UI elements.

\subsection{Prompt Engineering Tools and Scaffolding}

Several systems support prompt construction for LLMs.
PromptMaker~\cite{jiang2022promptmaker} provides a prototype-driven workflow
for NLP practitioners; PromptIDE and similar tools target developers.
PROPEL~\cite{strobelt2022interactive} offers an interactive prompt development
environment with output inspection.
For general users, tools like ChatGPT's ``custom instructions'' or
Anthropic's system-prompt editor expose partial scaffolding but require users
to already understand prompt structure.
Our work differs in two ways: we target \textit{novice} users of
\textit{local SLMs}, and we integrate scaffolding as a research
\textit{control}---not to help users write better prompts as an end goal,
but to hold prompting variance constant so that the SLM--LLM capability
comparison is interpretable.

\subsection{SLM Accessibility and Local Deployment}

Despite a growing ecosystem of quantized, edge-deployable models (Qwen,
Llama, Phi, Mistral), HCI research on SLM-specific user experience is sparse.
Studies of LLM-powered tools (e.g.~\cite{brachman2024enterprise}) characterize
knowledge-worker use patterns in cloud contexts; local deployment adds
constraints---model loading latency, context-length limits, output brittleness
---that change the user experience substantially.
Our Reddit study explicitly samples SLM-specific communities (\textit{r/LocalLLaMA},
\textit{r/SelfHosted}) and compares their task discourse with general LLM
communities (\textit{r/ChatGPT}, \textit{r/OpenAI}) to characterize where the
two use-practice spaces overlap.

\subsection{Comparative Evaluation of AI Models}

Prior work comparing models in user-facing settings typically varies only the
model while keeping the interface fixed~\cite{lee2022coauthor}, or uses
crowdsourced preference data (e.g., LMSYS Chatbot Arena~\cite{chiang2024chatbot}).
Neither approach controls for prompting-skill variance.
Easy\_SLM operationalizes a controlled within-subjects design in which the same
scaffolded input specification is submitted to both model classes, allowing
objective metric computation alongside user preference capture---a design we
believe is novel for the SLM--LLM comparison problem.

\subsection{Sustainable and Resilient Computing}

Crawford's CHI 2024 keynote calls for ``rematerializing AI''---attending to the
physical and energy costs of large model
inference~\cite{crawford2024rematerializing}.
Offline-first interfaces are increasingly relevant as infrastructure fragility
becomes visible~\cite{tully2025crowdstrike}.
Easy\_SLM contributes to this discourse by making the SLM/LLM trade-off
legible to users rather than leaving it implicit in deployment choices made
by technologists.

% =============================================================================
\section{Easy\_SLM System Design}
% =============================================================================

Easy\_SLM is a browser-based single-page application (Next.js, React,
TypeScript, Zustand) that routes the same user specification to either a local
SLM via the effGen API~\cite{effgen} or a cloud LLM (OpenAI / Anthropic) and
returns streamed responses through a unified Server-Sent Event contract.
The chat layout---colors, composer position, message bubbles---is identical
across conditions; the only difference is the backend.
This parity is deliberate: it prevents interface aesthetics from confounding
model capability comparisons (RQ3, RQ4).

We derived four design goals from the formative task-selection study
(Section~\ref{sec:tasksel}) and the prior literature on Gulf of
Envisioning~\cite{subramonyam2024gulf, zamfirescu2023johnny}.

% -----------------------------------------------------------------------------
\subsection{DG1: Paired SLM--LLM Execution with Objective Scoring}
% -----------------------------------------------------------------------------

\textit{Rationale.}
RQ1 and RQ2 require (i)~matched inputs submitted to both backends under
the same specification and (ii)~automatic metric computation.
Letting users freestyle prompts in separate tools confounds model capability
with prompting skill; a single shared interface eliminates that confound.

\textit{Implementation.}
The study home page presents three task cards, each with two entry buttons---
\textit{Begin with SLM} and \textit{Begin with LLM}---so participants explicitly
choose a condition per task.
A mode badge (CPU icon / Cloud icon) persists in the task bar throughout the
session.
Completion is tracked independently per task and mode: submitting the SLM
attempt does not mark the LLM attempt done, and vice versa.
The backend logs provider, model, task, and full prompt for post-hoc metric
computation (ROUGE, F1, rule-based pass rates) against human-written
references without exposing scores during the task.

% -----------------------------------------------------------------------------
\subsection{DG2: Translate User Intentions into Concrete Goals}
% -----------------------------------------------------------------------------

\textit{Rationale.}
Novice users tend to underspecify intent, and SLMs are especially brittle to
vague input~\cite{zamfirescu2023johnny}.
The \textit{intentionality gap}---not knowing what to ask---and the
\textit{capability gap}---not knowing what the model can produce---both
contribute to the Gulf of Envisioning~\cite{subramonyam2024gulf}.

\textit{Implementation.}
A two-step \textbf{Goal Wizard} modal guides users through (1)~writing a
one-sentence goal with an optional reframing prompt (e.g., ``turn my idea into
a clear, step-by-step task''), and (2)~selecting audience (\textit{just for me
/ team / general public}), output format (\textit{paragraph / checklist / table
/ JSON}), constraints, and success criteria.
In the SLM condition the wizard pre-loads a task-appropriate example goal that
participants can edit; in the LLM condition participants go directly to chat, so
the effect of goal scaffolding is isolatable.
The committed intent is compiled into a system prompt and displayed in a
persistent \textbf{Intent Ribbon} confirming the specification is live before
the first message is sent.

% -----------------------------------------------------------------------------
\subsection{DG3: Scaffold Prompt Design}
% -----------------------------------------------------------------------------

\textit{Rationale.}
Even after goal clarification, translating intent into effective instructions
remains hard---the \textit{instruction gap}.
Natural-language ambiguity and ``prompt brittleness'' cause SLM outputs to
vary unpredictably~\cite{zamfirescu2023johnny}, a confound that structured
decomposition can reduce.

\textit{Implementation.}
The \textbf{Prompt Plan Editor} lets users build an ordered list of prompt
steps.
\textit{Decompose with effGen} sends the current prompt text to effGen's
rule-based decomposition engine, returning subtasks that replace the step list
(no external LLM required for decomposition).
A \textbf{Variant Generator} produces three prompt variants (strict /
brief / explicit) side-by-side so users can inspect how phrasing changes
output before committing.
A \textbf{Template Gallery} offers four pre-built scaffolds---checklist,
summarize for audience, extract to JSON, rewrite with constraints---lowering
the entry floor for first-time users.
The compiled prompt is always visible in an accordion preview so users can
confirm what will be sent.

% -----------------------------------------------------------------------------
\subsection{DG4: Guided Model Selection}
% -----------------------------------------------------------------------------

\textit{Rationale.}
The SLM ecosystem offers hundreds of models varying in quantization, context
length, and hardware requirement.
Novice users face choice overload and lack heuristics for matching a model
to their task and device~\cite{subramonyam2024gulf}.

\textit{Implementation.}
The \textbf{Right Inspector} panel recommends models from a curated catalog
based on the committed intent and estimated compute tier.
A \textit{Reuse loaded model} toggle prevents costly model reloads across
turns.
In the LLM condition the same panel slot shows a provider selector
(OpenAI / Anthropic) and model dropdown, preserving layout parity while
removing SLM-specific guidance---so any right-panel effects measured in RQ4
are attributable to content, not position.

% =============================================================================
\section{Task Selection}
\label{sec:tasksel}
% =============================================================================

Choosing tasks that are both ecologically valid and shared between SLM and LLM
user communities is essential for meaningful capability comparison.
We conducted a formative Reddit study to ground task selection empirically
rather than relying on researcher intuition.

% -----------------------------------------------------------------------------
\subsection{Data Collection}
% -----------------------------------------------------------------------------

We scraped Reddit via the PullPush API over a twelve-month window
(March~2025--March~2026) using a two-stream keyword strategy designed to
capture distinct but potentially overlapping user populations.

\textbf{Stream~1} (SLM/local) targeted small-model discourse with 12 anchor
terms (e.g., \textit{local llm}, \textit{small language model},
\textit{edge ai}, \textit{self-hosted llm}) in the subreddits
\textit{r/LocalLLaMA}, \textit{r/LocalLLM}, and \textit{r/SelfHosted}.
\textbf{Stream~2} (LLM/general) targeted large-model discourse with 12 anchor
terms (e.g., \textit{chatgpt}, \textit{generative ai},
\textit{foundation model}) in \textit{r/ChatGPT}, \textit{r/OpenAI},
\textit{r/ClaudeAI}, \textit{r/LLM}, and \textit{r/Gemini}.
Both streams applied a post-retrieval substring filter of nine shared
use-practice terms (e.g., \textit{use case}, \textit{how do you use},
\textit{helps me}) to retain only records describing actual usage
behavior---not model discussion in the abstract.
This two-stage strategy required only \textbf{24 API queries} rather than the
$12 \times 9 = 108$ a cross-product design would need, while preserving the
same logical target (\textit{model-context} AND \textit{use-practice}).

After deduplication by Reddit ID the corpus contained
\textbf{3,292 records}: 1,149 from Stream~1 and 2,143 from Stream~2.

% -----------------------------------------------------------------------------
\subsection{LLM Annotation}
% -----------------------------------------------------------------------------

Each record was submitted to GPT-4o via the OpenAI Batch API (50\% lower cost
than synchronous calls; an incremental resume mechanism handled interrupted
batches).
The annotation prompt asked the model to identify all concrete daily tasks
mentioned in the text and to assign each to one of four themes---
\textit{Creation}, \textit{Information Search}, \textit{Advice},
\textit{Automation}---and ten sub-themes, following the enterprise-use taxonomy
of Brachman et al.~\cite{brachman2024enterprise}.
Per record the schema returned task descriptions, theme and sub-theme
assignments, and evidence spans; multi-task records stored fields
pipe-separated.

% -----------------------------------------------------------------------------
\subsection{Canonical Grouping and Overlap Analysis}
% -----------------------------------------------------------------------------

Free-text task labels were normalized in two passes.
First, 35 hand-crafted regex patterns collapsed surface variants into shared
labels (e.g., \texttt{summar(iz|is)(e|ing|ation)} $\to$
\textit{summarization}).
Second, sentence-transformer embeddings and agglomerative clustering were
applied to the 206 manually reviewed records to merge semantically similar
labels into canonical groups.

This produced \textbf{29 canonical groups} from 39 unique raw labels.
A task group was classified as \textit{shared} if at least one annotated
record per stream mapped to it; \textit{SLM-only} or \textit{LLM-only}
otherwise.
Of the 29 groups, \textbf{19 (65\%) were shared}, 8 were SLM-only, and 2
were LLM-only.
Coverage was 104/104 SLM records and 100/102 LLM records, confirming that
the shared task space accounts for the large majority of observed discourse.

Table~\ref{tab:overlap} lists the ten highest-frequency shared groups.
The high overlap---despite the communities using different model classes in
different infrastructure contexts---supports the premise that everyday task
needs are largely model-agnostic; the \textit{capability} to satisfy those
needs, not the tasks themselves, is what differs.

\begin{table}[t]
  \caption{Top shared canonical task groups ranked by combined mention count.}
  \label{tab:overlap}
  \small
  \begin{tabular}{lrrr}
    \toprule
    Canonical task & SLM & LLM & Total \\
    \midrule
    Creative writing      & 18 & 21 & 39 \\
    Summarization         & 18 & 11 & 29 \\
    Scripting             &  8 & 12 & 20 \\
    Task automation       &  8 & 11 & 19 \\
    Brainstorming         &  4 & 11 & 15 \\
    Transcription         & 11 &  2 & 13 \\
    Translation           &  5 &  3 &  8 \\
    Analyze data          &  2 &  6 &  8 \\
    Planning / scheduling &  2 &  4 &  6 \\
    Agent workflow        &  3 &  3 &  6 \\
    \bottomrule
  \end{tabular}
\end{table}

% -----------------------------------------------------------------------------
\subsection{Asymmetric Patterns and Design Implications}
% -----------------------------------------------------------------------------

Two asymmetries are notable.
\textit{Transcription} is SLM-dominant (11 vs.\ 2 mentions), consistent with
privacy concerns driving local audio processing---an SLM-specific use case
that warrants dedicated scaffolding support not required in the LLM condition.
\textit{Brainstorming} is LLM-dominant (4 vs.\ 11), suggesting that
open-ended ideation is currently under-served by local deployments; the
prompt-variant generator in DG3 directly targets this gap.
The 8 SLM-only groups (\textit{email management}, \textit{RAG},
\textit{knowledge base}, etc.) confirm that privacy-sensitive and
infrastructure-heavy workflows are a distinctive SLM use case beyond the
shared core.

% -----------------------------------------------------------------------------
\subsection{Final Task Shortlist}
% -----------------------------------------------------------------------------

For the comparative study we selected the highest-frequency shared task from
each of the four annotation themes, merging \textit{Advice} and
\textit{Automation} (low-count overlap) into a single planning condition:

\begin{itemize}
  \item \textbf{Task~1 -- Writing} (\textit{Creation} theme):
    \textit{creative writing} --- 18 SLM, 21 LLM mentions.
    Participants produce a practical short piece (email, social post, or
    short description) until the output is something they would ``realistically
    use or send.''
  \item \textbf{Task~2 -- Summarizing} (\textit{Information Search} theme):
    \textit{summarization} --- 18 SLM, 11 LLM mentions.
    Participants summarize a 250-word passage until the result is clear enough
    to ``use instead of re-reading the original.''
  \item \textbf{Task~3 -- Planning} (\textit{Advice} theme):
    \textit{planning / scheduling} --- 2 SLM, 4 LLM mentions.
    Participants produce a day-by-day plan for a given scenario until it feels
    ``realistic, specific, and easy to follow.''
\end{itemize}

The three tasks span a spectrum from open-ended creative generation (highest
LLM advantage expected) to structured output (where constraint scaffolding
may close the gap), enabling a graded test of the ``good enough'' boundary
(RQ2).

% =============================================================================
\section{Results and Discussion}
% =============================================================================

% -----------------------------------------------------------------------------
\subsection{RQ1 -- Task-Level Capability Gap}
% -----------------------------------------------------------------------------

[To be completed after data collection.
Table~\ref{tab:metrics} will report mean ROUGE-1/2/L, F1, and rule-based pass
rates by task and condition with Wilcoxon signed-rank $p$-values and effect
sizes ($r = Z/\sqrt{N}$), Bonferroni-corrected across three tasks.
Expected pattern: LLM outperforms SLM on ROUGE for open-ended writing, where
output space is large and reference coverage is low; the gap narrows for
summarization (well-defined reference) and further for planning (rule-based
checklist scoring that rewards structure over fluency).]

\begin{table}[t]
  \caption{Objective output quality by task and condition (mean $\pm$ SD).
           \textit{[To be replaced with study data.]}}
  \label{tab:metrics}
  \small
  \begin{tabular}{llcc}
    \toprule
    Task & Metric & SLM & LLM \\
    \midrule
    Writing    & ROUGE-L   & -- & -- \\
    Summarize  & ROUGE-L   & -- & -- \\
    Planning   & Pass rate & -- & -- \\
    \bottomrule
  \end{tabular}
\end{table}

% -----------------------------------------------------------------------------
\subsection{RQ2 -- ``Good Enough'' Boundary}
% -----------------------------------------------------------------------------

[To be completed.
``Good enough'' is operationalized as: SLM metric $\geq 90\%$ of LLM metric
AND participant preference not significantly favoring LLM ($p > .05$,
binomial test).
Hypothesis: summarization and planning will meet the threshold; open-ended
creative writing will not, reflecting the higher output variance of generation
tasks without a reference ceiling.
If confirmed, this boundary becomes an actionable decision rule: for
structured everyday tasks, a privacy-preserving local SLM is a justified
default, while fluency-sensitive generation warrants the higher-resource LLM.]

% -----------------------------------------------------------------------------
\subsection{RQ3 -- Prompting-Skill Confound}
% -----------------------------------------------------------------------------

[To be completed.
Prompt logs will be coded for constraint completeness, specificity, and format
adherence (two coders, Cohen's $\kappa$ target $> 0.7$).
If the Goal Wizard and Prompt Plan Editor reduce within-condition variance in
constraint completeness---compared to a free-form baseline---this confirms that
scaffolding acts as a \textit{methodological control}, not only as a user aid.
The implication is that uncontrolled SLM--LLM comparisons in the literature
likely overstate the capability gap by attributing prompting-skill variance to
the model.]

% -----------------------------------------------------------------------------
\subsection{RQ4 -- Interface Affordances}
% -----------------------------------------------------------------------------

[To be completed.
NASA-TLX scores will quantify the workload cost of scaffolding.
Format-violation rates (checklist items missing, JSON malformed, constraint
words absent) will test whether step decomposition reduces constraint errors.
Expected: SLM+scaffold shows higher task demand but fewer format violations
than SLM without scaffold; a non-significant workload difference between
LLM (no scaffold) and SLM+scaffold would support the argument that
scaffolding compensates for model brittleness at acceptable cognitive cost.]

% -----------------------------------------------------------------------------
\subsection{Discussion}
% -----------------------------------------------------------------------------

\subsubsection{Everyday task needs are largely model-agnostic.}
The 65\% overlap between SLM and LLM task discourse confirms that users bring
the same goals to both model classes; what varies is the capability to satisfy
them and the friction of doing so.
This reframes the research question from ``what do people use SLMs for?''
to ``which tasks can SLMs serve well enough?''---a more tractable and
actionable question.

\subsubsection{Scaffolding as research infrastructure.}
The Goal Wizard and Prompt Plan Editor are designed to serve dual roles: a
user aid that reduces the Gulf of Envisioning for novices, and a
\textit{variance-reduction mechanism} that makes SLM--LLM comparisons
interpretable.
Making this dual role explicit is a methodological contribution: future
comparative studies can adopt or adapt Easy\_SLM's scaffolding layer as
a controlled prompting substrate, and the interface is open-sourced for reuse.

\subsubsection{Sustainable trade-offs made legible.}
The ``good enough'' boundary operationalized in RQ2 translates a technical
model-selection decision into user-facing information.
Displaying task-level sufficiency evidence in the interface---rather than
hiding it in deployment infrastructure choices---aligns with Crawford's call
to rematerialize AI trade-offs~\cite{crawford2024rematerializing}.
Users in the study who understand that SLM is ``good enough'' for their
planning task but not for their creative writing task are better equipped to
make intentional, sustainable choices about cloud dependency.

\subsubsection{Limitations.}
The formative Reddit study over-represents technically engaged users;
casual SLM adopters may report different task patterns.
ROUGE and rule-based pass rates are imperfect proxies for utility on
open-ended tasks and should be interpreted alongside preference and confidence
data.
Results are specific to the evaluated SLM [\textit{model name}, quantization];
different models may shift the capability gap substantially.
The within-subjects design introduces potential demand characteristics:
participants may moderate their behavior once they understand the comparison
purpose.

\subsubsection{Conclusion.}
Easy\_SLM pairs a formative task-discovery pipeline with a scaffolded
comparison interface to produce empirically grounded, confound-controlled
evidence about when SLMs are sufficient for everyday use.
The 19 shared tasks identified across SLM and LLM communities provide a
reusable task taxonomy for future work; the four design goals provide a
template for fair capability comparison interfaces; and the ``good enough''
framework offers a decision procedure for sustainable model selection that
is actionable for everyday users rather than exclusively for researchers
or system administrators.


\begin{acks}
[To be completed before submission.]
\end{acks}

\bibliographystyle{ACM-Reference-Format}
\bibliography{references}

\end{document}
